\ListokName{18. Вероятность}

\shapkaUstnoA


Бросим симметричную монету тысячу раз. Предсказать, как она упадёт в очередной раз, нельзя, но орлов, скорее всего, будет примерно половина. Бросок монеты~— \textit{эксперимент}: процедура, которую можно повторять сколько угодно раз, а результат от раза к разу разный. \textit{Событие}~— утверждение о результате эксперимента: после каждого опыта можно сказать, произошло оно или нет: «выпал орёл», «на кубике выпало чётное число». Если в $N$ повторениях эксперимента событие произошло $k$ раз, число $k/N$ называют \textit{частотой} события.

У симметричного кубика грани ничем не отличаются друг от друга, поэтому в длинной серии бросков каждая грань обычно выпадает примерно в шестой части случаев. Слова «случайный выбор» («тянут наугад», «садится на случайное место») означают то же самое: в длинной серии опытов каждый вариант выбора встречается примерно одинаково часто.

\zp Симметричный кубик бросают много раз. Какова примерная частота выпадения \lett шестёрки; \leth чётного числа; \leth числа, делящегося на 3; \lett числа, не делящегося ни на 2, ни на 3?

\zn Каждый урок геометрии Наталия Павловна случайно выбирает одного из учеников вашего класса и просит стереть с доски. С какой частотой  Андрей  стирает  с доски на  уроке геометрии? \lett Что  влияет на эту частоту (кроме обыкновения Наталии Павловны просить одних учеников стирать с доски чаще, чем других)?

\zn В пруду обитают караси двух видов. Выловили 40 карасей. Оказалось, что 12 из них первого вида. Какой прогноз о доле карасей первого вида среди обитающих в пруду карасей вы можете сделать?
\lett Ихтиолог Сева выловил $k$ карасей, прикрепил к их плавнику метку, а потом выпустил их обратно в пруд. На следующий день он выловил $N$  карасей и $n$ из них оказались с меткой. Какой вывод можно сделать в результате этого исследования?


\bigskip

Если эксперимент повторять много раз, частота события колеблется около одного и того же числа~— где бы и кем бы эксперимент ни проводился. Это число называют \textit{вероятностью} события $A$ и обозначают $P(A)$. Это не теорема, а наблюдение о мире; на нём основана вера в то, что теорию вероятностей можно применять к реальной жизни. Слово «вероятность» используют и в других смыслах: например, говорят о вероятности единичного события, которое нельзя повторить много раз,~— что определённая команда выиграет завтрашний матч. В этом листке мы говорим только о вероятности как о частоте.

Событие «произошло хотя бы одно из событий $A$ и $B$» обозначают $A\cup B$, «произошли и $A$, и $B$»~— $A\cap B$, «$A$ не произошло»~— $\overline{A}$; его называют \textit{противоположным} к $A$. События $A$ и $B$ \textit{несовместны}, если они не могут произойти одновременно.




\z Монету бросили 10 раз. Какова вероятность того, что 
\lett 10 раз выпал орёл;
\lett сначала выпало 5 орлов, а затем 5 решек;
\lett выпало 5 орлов и 5 решек (в любом порядке)?


\z Подбрасывают два кубика --- синий и зелёный. С какой вероятностью получится, что
\lett на синем выпало чётное число, а на зелёном 1 или 6?
\lett сумма очков на двух кубиках равна 5?


\z В шляпе лежат карточки со словами: 10 с надписью  <<дуб>>, 5 с надписью <<зуб>>,  15 с надписью <<куб>> и 20 с надписью <<зубр>>. Вытаскивают одну случайную карточку. Рассмотрим события:\\ $A$=<<слово заканчивается на букву ``б''>>\\ $B$=<<слово начинается на букву ``з''>>\\$C$=<<в слове есть буква ``у''>>\\ $D$=<<слово начинается на букву ``б''>>. Верны ли утверждения:
\lett $P(C)+P(D)=P(C\cup D)$;
\leth $P(A)+P(B)=P(A\cup B)$.

\newpage

\zncirc  Объясните, почему $P(A\cup B)=P(A)+P(B)$, если события $A$ и $B$ несовместны. \leth Чему равна вероятность события, которое происходит при любом результате эксперимента? \lett Объясните, почему $P(A)+P(\overline{A})=1$.

\zpcirc Верно ли, что $P(A\cup B)=P(A)+P(B)$ для любых событий $A$ и $B$? Если нет, сформулируйте верное равенство.


\zp Вероятность события $A$ равна 1/3, вероятность события $B$ равна 1/2, вероятность события «произошло событие $A$ или событие $B$» равна 2/3. Как такое возможно? Придумайте пример эксперимента и событий, для которых будут такие вероятности.


\zp В классе 25 детей. Для дежурства наугад выбирают двоих. Вероятность того, что оба дежурных окажутся мальчиками, равна 3/25. Сколько в классе девочек?


\vspace{-.2cm}

$$\star\star\star$$

\textit{Вероятностная модель} эксперимента~— это конечное множество $\Omega$ и числа $p(\omega)\ge 0$, по одному для каждого $\omega\in\Omega$, в сумме равные 1. Элементы $\Omega$~— \textit{элементарные исходы}: несовместные события, из которых при каждом проведении эксперимента происходит ровно одно; $p(\omega)$~— их вероятности. Событию отвечает подмножество $A\subset\Omega$, и $P(A)$~— сумма $p(\omega)$ по всем $\omega\in A$. Один эксперимент можно описать разными моделями; если обе правильно описывают эксперимент, то вероятность события в обеих моделях одна и та же.

\zpdag Докажите, исходя из определения модели: \textbf{а)} $P(\Omega)=1$ и $P(\emptyset)=0$; \textbf{б)} если $A\subset B$, то $P(A)\le P(B)$; \textbf{в)} $P(\overline{A})=1-P(A)$; \textbf{г)} $P(A\cup B)=P(A)+P(B)-P(A\cap B)$.

\zp Каждый ученик подбрасывает один кубик, одну монетку и вытаскивает из мешка с полным набором домино одну из доминошек. Постройте вероятностную модель этого эксперимента. Что будет элементарным исходами? Сколько элементарных исходов в этой модели? Чему равны их вероятности? Приведите пример пары несовместных событий в этой модели, а также пример пары противоположных событий.

\zp С какой вероятностью в предыдущей задаче выпадет орёл? Можно ли было   для ответа на этот вопрос  построить другую вероятностную модель?

\z Почтальон Печкин бросает монетку. Если выпал орёл, бросает монетку ещё два раза. Если оба раза выпадет решка, идёт на работу пешком, а иначе едет  на велосипеде. Если выпадает решка, бросает кубик. Если выпадает 5 или 6, идёт пешком, иначе едет на автобусе. Если Печкин едет на велосипеде, то с вероятностью 70\%  успевает развезти все письма во время, а если идёт пешком --- то с вероятностью~40\%. \lett Постройте вероятностную модель.
\lett Придумайте какое-нибудь событие для этого эксперимента и предложите однокласснику найти его вероятность.
\lett Найдите вероятность события в этом эксперименте, придуманного кем-нибудь из ваших одноклассников.
